\section{A logarithmic oracle lower bound}
\label{sec:lower}

The lower bound already holds for centered Gaussian targets, whose
minimizer is known. The obstruction is therefore present in the same
oracle model as the upper bound. We adapt the Gaussian block-Krylov
strategy of
\citet[Section~5.3.3, Lemma~39 and Theorem~45]{chewi2023lower},
which already yields logarithmic dependence on dimension at fixed
condition number. We give a complete argument for the spectral interval
$[1,2]$ and total-variation accuracy $1/10$ used here, covering arbitrary
measurable adaptive algorithms and unrestricted real computation.

\begin{theorem}\label{thm:lower}
There are absolute constants $c>0$ and $d_0$ such that $Q(d)\ge c\log d$
for every $d\ge d_0$. This bound holds even when the inputs are restricted
to potentials $V(x)=x^TAx/2$ with $I\preceq A\preceq2I$.
The query cap is imposed on every execution, and the algorithm may use
arbitrary private randomness independent of its input.
\end{theorem}

For these potentials the gradient reply is $Ax$, and the value
$x^TAx/2$ is already determined by the query and its gradient reply.
Thus an exact value-and-gradient query is equivalent to one query to
$x\mapsto Ax$. Each potential belongs to $\mathcal V_d$, has minimizer
$0$, and has target $N(0,A^{-1})$.

The proof has three steps. First, deferred revelation of a random
rotation represents any output in a Gaussian block-Krylov span, including
outputs outside the span of the actual queries and replies. Second, a
cosine statistic is small uniformly on this enlarged random span.
Finally, the same statistic has a large bias under the target Gaussian,
which separates the two joint experiments in total variation. These
steps use no result from the sampling upper bound.

\subsection{Adaptive queries and unrevealed rotations}

We first record a representation of an arbitrary algorithm's output.
The auxiliary Gaussian vectors in this representation are independent of
one another, but need not be independent of the input matrix.

\begin{lemma}[Deferred rotation]\label{lem:deferred-haar}
Let $\Lambda$ be a fixed real diagonal $d\times d$ matrix, let $O$ be
Haar distributed on the orthogonal group, and set $A=O\Lambda O^T$.
Consider any measurable randomized algorithm, with private randomness
independent of $O$, making at most $q$ adaptive queries to the map $x\mapsto Ax$, where $2q+1<d$, and producing an output
$X\in\R^d$. There is an extension of this experiment carrying independent
standard Gaussian vectors $g_1,\ldots,g_{q+1}\in\R^d$ such that
\begin{equation}\label{eq:krylov-span}
 O^TX\in\operatorname{span}
 \{\Lambda^j g_i:1\le i\le q+1,\ 0\le j\le q\}
 \qquad\text{almost surely}.
\end{equation}
The joint law of the original variables, including the algorithm's
randomness and $O$, is unchanged. The vectors $g_i$ need not be
independent of $O$ or of $A$.
\end{lemma}

\begin{proof}
Include all the algorithm's randomness, including any randomness used at
output, in a private seed independent of $O$. Fix this seed for the
moment. Pad executions with zero queries so that there are exactly $q$
query slots.

We use the following conditional property of Haar measure. For a fixed
restriction $T:E\to W$ of an orthogonal map, the conditional law of its
remaining restriction is uniform among isometries from $E^\perp$ to
$W^\perp$. For a specified unit vector in either complement, its partner
is uniform on the unit sphere of the other complement. After revealing
that pair, the map between the smaller complements is again uniform.
This follows from invariance under the orthogonal transformations fixing
$E$ and $W$ pointwise. Applying this rule successively also permits each
specified direction to be a measurable function of the already revealed
data.

All conditional laws below are laws of finite-dimensional data after the
seed has been fixed. To represent the subspaces and partial isometries
measurably, apply Gram--Schmidt to the revealed vectors in their order of
appearance, discard zero residuals, and complete a basis by scanning the
standard coordinate vectors in order. These operations are Borel on
each rank stratum. Thus regular conditional laws and the successive
Gaussian draws can be chosen as measurable kernels; no continuity of
the algorithm's query maps is required.

Initially $E=W=\{0\}$. After each reply, reveal the restriction
$T=O|_E$, where $E$ is the span of the eigenframe coordinates of all
queries and replies so far, and $W=OE$ is their span in the original
coordinates. Our conditional invariant is that the unrevealed map is
Haar on the two complements, given the full latent history: the seed,
the transcript, this restriction, and all auxiliary variables constructed
so far.

Let $x$ be the next query, which is fixed given that history. Write
$x=x_W+x_\perp$ according to $W\oplus W^\perp$. The preimage of $x_W$
is known. If $x_\perp\ne0$, put
\[
 w=\frac{x_\perp}{|x_\perp|},\qquad v=O^Tw\in E^\perp.
\]
Conditionally on the history, $v$ is uniform on the complementary unit
sphere. Let $h$ be a standard Gaussian vector in $E$, and let $a$ have
the distribution of the Euclidean norm of a standard Gaussian in
$\R^{d-\dim E}$. Draw $h$ and $a$ independently of one another and of
the unrevealed map, conditional on the history, and define
\[
                         g_i=h+av.
\]
The Gaussian polar decomposition shows that $g_i$ is a standard Gaussian
in $\R^d$ conditional on the entire previous history. Moreover, $a>0$ almost surely, so
$v=(g_i-h)/a\in E+\operatorname{span}\{g_i\}$. Reveal $Ov=w$ and extend $T$ to
$E+\operatorname{span}\{v\}$. Conditioning on the full vector $g_i$
does not reveal any further part of $O$: its parallel component $h$
and complementary radius $a$ are independent auxiliary variables, and
its normalized complementary component is precisely $v$. Thus the
conditional Haar invariant persists. If $x_\perp=0$, instead draw a fresh
independent standard Gaussian $g_i$ and make no extension at this step.

In either case the query preimage $u=O^Tx$ is now known and lies in the
old space $E+\operatorname{span}\{g_i\}$. Its reply has preimage
$\Lambda u$. Project $\Lambda u$ onto the enlarged eigenframe space.
The image of this projection is known. If the residual is nonzero,
reveal the image under $O$ of its unit direction and extend the
isometry once more. This is the other form of the same Haar-revelation
rule: the direction in the eigenframe is fixed by the current history,
and its image is uniform in the unrevealed complementary sphere. The
resulting reply is exactly $Ax$, and the conditional invariant is
preserved.

Every $g_i$ has conditional law $N(0,I_d)$ given all earlier Gaussian
vectors and the earlier history. Consequently these vectors are
independent. To check the polynomial degrees, let $E_i$ be the revealed
eigenframe space after reply $i$. Induction gives
\[
 E_i\subseteq\operatorname{span}
 \{\Lambda^j g_a:1\le a\le i,\ 0\le j\le i\}.
\]
Indeed, the new query preimage is in $E_{i-1}+\operatorname{span}\{g_i\}$,
and adjoining its reply applies $\Lambda$ once. All old directions
remain in the indicated larger span.

Finally treat the algorithm's output as one more query direction, with
no reply: repeat the preimage construction using $g_{q+1}$. This gives
$O^TX\in E_q+\operatorname{span}\{g_{q+1}\}$, proving
\eqref{eq:krylov-span}. At most two dimensions were exposed per oracle
query and one for the output, so the assumed dimension inequality
ensures that every required complement is available. The construction
started with the genuine Haar input and never altered it. Its added
variables are obtained by measurable kernels from the original variables
and independent auxiliary draws. The construction therefore applies
jointly measurably for the fixed-seed algorithms; integrating over the
original independent seed proves the assertion for randomized algorithms.
\end{proof}

\subsection{A statistic invisible on the block span}

For the rest of the section define
\begin{equation}\label{eq:lower-spectrum}
 \theta_k=\frac{(k-1/2)\pi}{d},\qquad
 \lambda_k=\frac{3+\cos\theta_k}{2},\qquad
 \Lambda=\operatorname{diag}(\lambda_1,\ldots,\lambda_d),
 \quad 1\le k\le d.
\end{equation}
All eigenvalues lie strictly between $1$ and $2$. For an integer $q\ge0$
put
\[
 n=2q+1,\qquad p_k=\cos(n\theta_k),\qquad
 D_p=\operatorname{diag}(p_1,\ldots,p_d).
\]
The frequency $n$ exceeds twice the polynomial degree in
\eqref{eq:krylov-span}. This makes the quadratic form $u^TD_pu$ small
simultaneously for all vectors in that random span.

\begin{lemma}\label{lem:lower-span-statistic}
Suppose $4q+1<2d$, and let $g_1,\ldots,g_{q+1}$ be independent standard
Gaussian vectors. Set
\[
 \mathcal K_q=\operatorname{span}
 \{\Lambda^j g_i:1\le i\le q+1,\ 0\le j\le q\},\qquad
 \tau_d=\frac{40(q+1)^2}{\sqrt d}.
\]
If $\tau_d\le1/2$, then, with probability at least $0.99$,
\begin{equation}\label{eq:lower-uniform-statistic}
 |u^TD_pu|\le2\tau_d|u|^2
 \qquad\text{for every }u\in\mathcal K_q.
\end{equation}
\end{lemma}

\begin{proof}
Let $m=(q+1)^2$ and form a $d\times m$ matrix $B$, whose columns are
indexed by $(i,j)$ with $1\le i\le q+1$ and $0\le j\le q$, by
\[
 B_{k,ij}=d^{-1/2}(g_i)_k c_j\cos(j\theta_k),
 \qquad c_0=1,\quad c_j=\sqrt2\quad(j\ge1).
\]
The functions $\cos(j\theta)$ are degree-$j$ polynomials in
$\cos\theta$ with nonzero leading coefficients. Since
$\cos\theta_k=2\lambda_k-3$, the range of $B$ is precisely
$\mathcal K_q$.

Discrete cosine orthogonality gives
\begin{equation}\label{eq:lower-gram-means}
                  \E B^TB=I_m,\qquad \E B^TD_pB=0.
\end{equation}
For completeness, the equally spaced half-grid satisfies
\begin{equation}\label{eq:lower-cosine-sum}
 \sum_{k=1}^d\cos(\ell\theta_k)
 =\frac{\sin(\ell\pi)}{2\sin(\ell\pi/(2d))}=0,
 \qquad 1\le\ell<2d,\quad \ell\in\mathbb Z.
\end{equation}
Here the finite cosine sum follows by summing a geometric progression,
and its denominator is nonzero in the indicated range.
Expand products of two cosines for the first identity in
\eqref{eq:lower-gram-means}. For the second, expand the product of
three cosines with frequencies $n,j,j'$. Its frequencies are nonzero
because $n>j+j'$, and have absolute value at most $4q+1<2d$.
Different Gaussian blocks have zero cross means in both cases.

An entry of either Gram matrix is a sum of independent row terms.
The deterministic coefficient in each term has absolute value at most
$2/d$, and the variance of a product of two standard Gaussian entries
is at most $2$. Thus every entry has variance at most $8/d$, and
\[
 \E\bigl[\|B^TB-I_m\|_{\HS}^2+
               \|B^TD_pB\|_{\HS}^2\bigr]
 \le\frac{16m^2}{d}.
\]
By Markov's inequality,
\[
 \P\!\left(\|B^TB-I_m\|_{\HS}^2+
                 \|B^TD_pB\|_{\HS}^2>\tau_d^2\right)
 \le\frac{16m^2/d}{1600m^2/d}=0.01.
\]
Since $\|C\|_{\op}\le\|C\|_{\HS}$, outside this event both operator
norms are at most $40m/\sqrt d=\tau_d$. Whenever these bounds hold, for every
$a\in\R^m$,
\[
 |Ba|^2\ge(1-\tau_d)|a|^2\ge\tfrac12|a|^2,
 \qquad |(Ba)^TD_p(Ba)|\le\tau_d|a|^2.
\]
These inequalities imply \eqref{eq:lower-uniform-statistic} on the
entire range of $B$.
\end{proof}

The uniformity in this lemma is essential: the output's coefficients in
the block span can depend arbitrarily on the whole transcript. No
independence between these coefficients, the input rotation, and the
Gaussian vectors is required.

\subsection{Separation from the target Gaussian}

An ideal target has a detectable bias at the same frequency. Write
\[
                         \varrho=3-\sqrt8\in(0,1).
\]
The geometric-series formula gives the absolutely convergent identity
\begin{equation}\label{eq:lower-fourier}
 \frac{2}{3+\cos\theta}
 =\frac1{\sqrt2}
   \left(1+2\sum_{j\ge1}(-\varrho)^j\cos(j\theta)\right).
\end{equation}
Indeed, the series in parentheses sums to
$(1-\varrho^2)/(1+2\varrho\cos\theta+\varrho^2)$, and
$1+\varrho^2=6\varrho$ and $1-\varrho^2=4\sqrt2\,\varrho$.

\begin{lemma}\label{lem:lower-target-statistic}
Let $U\sim N(0,\Lambda^{-1})$, with $\Lambda$ from
\eqref{eq:lower-spectrum}, and suppose $n=2q+1<d$. Then
\begin{align}
 \mu_{d,q}:=\E[U^TD_pU]
 &=\frac d{\sqrt2}(-\varrho)^n+R_{d,q},
 & |R_{d,q}|&\le
   \frac{\sqrt2\,d\,\varrho^{2d-n}}{1-\varrho},
 \label{eq:lower-target-mean}\\
 \Var(U^TD_pU)&\le2d,
 & \E|U|^2&\le d,\qquad \Var(|U|^2)\le2d.
 \label{eq:lower-target-variance}
\end{align}
\end{lemma}

\begin{proof}
Independence of the coordinates gives
$\mu_{d,q}=\sum_k p_k/\lambda_k$. Substitute
\eqref{eq:lower-fourier}. The constant term vanishes by
\eqref{eq:lower-cosine-sum}. For $1\le j<2d-n$, use
\[
 \cos(n\theta_k)\cos(j\theta_k)
 =\tfrac12\cos((n-j)\theta_k)+\tfrac12\cos((n+j)\theta_k).
\]
If $j\ne n$, both frequencies have nonzero absolute value below $2d$,
so their sums vanish. If $j=n$, the constant term contributes $d/2$
and the frequency $2n<2d$ contributes zero. Thus only $j=n$ contributes
in this range. Each remaining cosine product has
absolute sum at most $d$, so summing the geometric tail proves
\eqref{eq:lower-target-mean}. Finally, a centered normal variable with
variance $v$ satisfies $\Var(U_k^2)=2v^2$. Since
$|p_k|\le1$ and $\lambda_k\ge1$, summing these coordinate variances
proves \eqref{eq:lower-target-variance}.
\end{proof}

\begin{proof}[Proof of Theorem~\ref{thm:lower}]
Take $A=O\Lambda O^T$ with $O$ Haar and $\Lambda$ as in
\eqref{eq:lower-spectrum}. The reduction at the start of the section identifies an exact
value-and-gradient algorithm on $V_A(x)=x^TAx/2$ with an algorithm
querying $x\mapsto Ax$, with the same query cap and output law.

Fix any algorithm using at most $q$ queries on every execution, where
\[
 0\le q\le c_0\log d,\qquad
 c_0=\frac1{16\log(1/\varrho)}.
\]
All dimension restrictions in Lemmas~\ref{lem:deferred-haar},
\ref{lem:lower-span-statistic}, and \ref{lem:lower-target-statistic}
hold for sufficiently large $d$, uniformly over these $q$.
Let $\P_{\mathrm{alg}}$ be the joint law of the Haar rotation and the
algorithm's output on this input. Define the event on pairs $(O,x)$
\begin{equation}\label{eq:lower-separating-event}
 \mathcal E_{d,q}=
 \left\{(O,x): |O^Tx|^2\le2d,\quad
       \left|(O^Tx)^TD_p(O^Tx)\right|>
                         160(q+1)^2\sqrt d\right\}.
\end{equation}
This event is measurable. It is also determined by $(A,x)$: if $T_n$
is the polynomial defined by $T_n(\cos\theta)=\cos(n\theta)$, then
$OD_pO^T=T_n(2A-3I)$. We retain $O$ in the notation to use eigenframe
coordinates throughout the comparison.

Apply Lemma~\ref{lem:deferred-haar} to the algorithm's output and then
Lemma~\ref{lem:lower-span-statistic} to its auxiliary Gaussian vectors.
On the latter lemma's event, the norm condition in
\eqref{eq:lower-separating-event} forces the quadratic form to be at
most $2\tau_d(2d)=160(q+1)^2\sqrt d$. Hence
\begin{equation}\label{eq:lower-algorithm-event}
                       \P_{\mathrm{alg}}(\mathcal E_{d,q})\le0.01.
\end{equation}

For the target joint law $\P_{\mathrm{target}}$, sample $O$ from the same
prior and then $X$ from $N(0,O\Lambda^{-1}O^T)$. Conditional on $O$, the vector
$U=O^TX$ has law $N(0,\Lambda^{-1})$. Our choice of $c_0$ gives
\[
 \varrho^{2q+1}\ge\varrho d^{-1/8},\qquad
 \frac{d\varrho^{2q+1}}{(q+1)^2\sqrt d}
 \ge\frac{\varrho d^{3/8}}{(1+c_0\log d)^2}\longrightarrow\infty.
\]
The remainder in \eqref{eq:lower-target-mean} satisfies
\[
 \frac{|R_{d,q}|}{d\varrho^n/\sqrt2}
 \le\frac{2\varrho^{2d-2n}}{1-\varrho}\longrightarrow0
\]
uniformly, since $n\le1+2c_0\log d$. Thus it is negligible relative
to the leading term. In particular,
with $T_{d,q}=160(q+1)^2\sqrt d$, we have
$|\mu_{d,q}|>T_{d,q}$ for large $d$, and
\[
 \P_{\mathrm{target}}(\mathcal E_{d,q}^c)
 \le\frac2d+
       \frac{2d}{(|\mu_{d,q}|-T_{d,q})^2}=o(1)
\]
uniformly, by Lemma~\ref{lem:lower-target-statistic} and Chebyshev's
inequality. Thus the probabilities of $\mathcal E_{d,q}$ under the
two joint experiments differ by at least $0.99-o(1)$.

To compare these joint laws with the required guarantee on each fixed
input, let $K_O$ denote the algorithm's output law on $A=O\Lambda O^T$
and let $G_O=N(0,A^{-1})$. For a measurable joint event $\mathcal E$,
write $\mathcal E_O=\{x:(O,x)\in\mathcal E\}$. Both laws use the same
Haar prior, so
\begin{align*}
 \left|\P_{\mathrm{alg}}(\mathcal E)
       -\P_{\mathrm{target}}(\mathcal E)\right|
 &=\left|\int\bigl[K_O(\mathcal E_O)-G_O(\mathcal E_O)\bigr]
                   \,d\operatorname{Haar}(O)\right|\\
 &\le\int\|K_O-G_O\|_{\TV}\,d\operatorname{Haar}(O).
\end{align*}
The pointwise sampling guarantee would bound the last integral by
$1/10$. The event used to disprove that guarantee is a test of the joint
laws; the algorithm need not know or compute it.
This is impossible for \eqref{eq:lower-separating-event}. Therefore no
such algorithm with $q\le c_0\log d$ can meet the required accuracy for
all inputs. Decreasing the constant if necessary proves the theorem.
\end{proof}
